\section{De la presentación de lanzamiento al desempeño real: cambios en la visibilidad de las empresas de modelos}

En la segunda mitad, Zhang Peng (张鹏) menciona un cambio muy sutil. Antes, el lanzamiento de un modelo solía requerir una presentación solemne, cientos de asistentes y una fuerte campaña de difusión. En cambio, tras la liberación en código abierto de GLM-4.5, Zhipu (智谱) no hizo un lanzamiento especialmente grande y, aun así, desarrolladores y productos del extranjero lo valoraron muy bien. Incluso hubo productos como Windsurf que sustituyeron su modelo o integraron este, y otros fabricantes que lo destilaron o recortaron. Este cambio indica que la visibilidad técnica de las empresas de modelos grandes está pasando de la «narrativa de la presentación de lanzamiento» al «desempeño en el uso real».

\subsection{Por qué disminuye el rendimiento marginal de las presentaciones de lanzamiento}

Antes se dijo que la visibilidad de un modelo se desplaza hacia su desempeño real; esta sección explica por qué. Cuando se lanzan demasiados modelos, la narrativa deja de ser escasa; lo escaso pasa a ser el uso real y la eficiencia en costos.

Después de varias rondas de lanzamientos en la industria de los modelos grandes, usuarios y desarrolladores ya se cansaron de «otro lanzamiento de modelo más». La publicidad puede atraer atención a corto plazo, pero que un modelo llegue a integrarse en el trabajo de los desarrolladores, que los clientes paguen por él y que la comunidad lo discuta por iniciativa propia depende cada vez más de su capacidad real y de su eficiencia en costos. La influencia de DeepSeek también radica en que llevó al primer plano de la industria el camino de «hablar poco, liberar mucho en código abierto y dejar que los resultados hablen».

\begin{knowledgebox}{Tres fuentes de visibilidad de un modelo}
\begin{tabularx}{\textwidth}{p{0.22\textwidth}p{0.34\textwidth}X}
\toprule
Fuente & Ventajas & Riesgos \\
\midrule
Presentación de lanzamiento & Narrativa concentrada; facilita crear hitos & Cansancio de los usuarios; difícil demostrar valor a largo plazo. \\
Benchmark & Comparable, difundible, útil para el juicio técnico & Propenso al sobreajuste a las clasificaciones; distante de las tareas reales. \\
Uso real & Retroalimentación natural de desarrolladores, clientes y productos & Retroalimentación lenta e incontrolable, pero la más cercana al valor comercial. \\
\bottomrule
\end{tabularx}
\end{knowledgebox}

\subsection{Ser usado también es una forma de evaluación}

Cuando los desarrolladores integran un modelo, cuando un producto lo usa para reemplazar a su modelo anterior, o cuando otros equipos lo destilan o lo reempaquetan, deja de ser solo un modelo de un paper o de una clasificación y entra en un ecosistema real. Esto tiene ventajas y también complicaciones: la ventaja es que tareas reales validan la capacidad del modelo; la complicación es que la captura de valor se vuelve más compleja y la influencia del código abierto no necesariamente se convierte de forma directa en ingresos.

\begin{importantbox}{La lógica de evaluación del desempeño real}
Una empresa de modelos no puede preguntarse solo «¿en qué lugar estamos en la clasificación?»; también debe preguntarse «¿quién lo usa, por qué lo usa, a quién reemplazó, si el costo es menor y si los clientes están dispuestos a seguir pagando?». Esa es la demostración de capacidad más importante en la etapa de empresa que cotiza en bolsa.
\end{importantbox}

\subsection{Resumen de la sección}

El centro de los lanzamientos de modelos se está desplazando de «yo digo que soy muy bueno» a «qué lograron otros usando mi modelo». Esta es también la nueva relación entre el ecosistema de código abierto, la comercialización y la capacidad de los modelos.
